\chapter{Scalar and Predictive Quantization}

\section{Table of Contents}

\begin{itemize}
  \item Core Idea
  \item Key Concepts
  \item Theory and Formulas
  \item Visual Schemes
  \item Examples
\end{itemize}

\section{Core Idea}

\textbf{Quantization} maps a continuous or high-precision signal to a finite set of reconstruction values. It is the main irreversible step in lossy compression: it reduces rate, but introduces distortion.

Scalar quantization alone is weak for natural images and audio because it treats each sample independently. \textbf{Predictive quantization} first removes predictable content, then quantizes the prediction error. If the prediction error has much lower variance and is entropy-coded, the rate-distortion gain is large.

\begin{calloutbox}{red!70!black}{Central message}
Quantization controls distortion; prediction reduces variance before quantization; entropy coding converts sparse prediction errors into real bitrate savings.
\end{calloutbox}

\section{Key Concepts}

\subsection{Scalar Quantizer Model}

A scalar quantizer $Q$ is defined by:

\begin{itemize}
  \item \textbf{Thresholds}:
\end{itemize}

\[
-\infty=t^0<t^1<\cdots<t^{L-1}<t^L=+\infty
\]

\begin{itemize}
  \item \textbf{Quantization regions}:
\end{itemize}

\[
\Theta^i=[t^{i-1},t^i), \qquad i=1,\ldots,L
\]

\begin{itemize}
  \item \textbf{Reconstruction levels}:
\end{itemize}

\[
\hat{x}^1,\hat{x}^2,\ldots,\hat{x}^L
\]

Mapping rule:

\[
x\in\Theta^i \Rightarrow Q(x)=\hat{x}^i
\]

Key quantities:

\begin{center}\fitwidth{%
\begin{tabular}{|l|l|l|}
\hline
\textbf{Quantity} & \textbf{Formula} & \textbf{Meaning} \\ \hline
Quantization error & $e=X-Q(X)=X-\tilde{X}$ & sample error \\ \hline
Distortion & $D=\sigma_Q^2=\mathbb{E}[(X-Q(X))^2]$ & MSE due to quantization \\ \hline
Rate & $R=\log_2 L$ & fixed-length bits/sample \\ \hline
SNR & $10\log_{10}(\sigma_X^2/\sigma_Q^2)$ & signal quality in dB \\ \hline
\end{tabular}}
\end{center}

\subsection{Uniform Quantization}

\textbf{Uniform quantization (UQ)} uses equal-width cells inside dynamic range $[-A/2,A/2]$:

\[
\Delta=\frac{A}{L}=A2^{-R}
\]

Reconstruction levels are cell centers:

\[
\hat{x}^i=-\frac{A}{2}+\left(i-\frac{1}{2}\right)\Delta
\]

For a midrise quantizer:

\[
Q(x)=\Delta\left\lfloor\frac{x}{\Delta}+\frac{1}{2}\right\rfloor
\]

with saturation at $\pm A/2$.

\begin{center}\fitwidth{%
\begin{tabular}{|l|l|l|l|}
\hline
\textbf{Type} & \textbf{Number of levels} & \textbf{Zero reconstruction level} & \textbf{Typical property} \\ \hline
\textbf{Midtread} & odd & yes & dead-zone around zero \\ \hline
\textbf{Midrise} & even & no & common for signed fixed-rate quantization \\ \hline
\end{tabular}}
\end{center}

\subsection{Dead-Zone}

A \textbf{midtread} quantizer keeps zero as a reconstruction level, so the central cell around zero behaves as a \textbf{dead-zone}: any sample with small magnitude falls inside it and maps exactly to zero. A \textbf{dead-zone quantizer} widens this zero region by a factor $\tau$, so even more low-amplitude samples collapse to zero.

Index (forward) and reconstruction (inverse), extending the midtread zero region:

\[
i=
\begin{cases}
\operatorname{sign}(x)\left\lfloor\dfrac{|x|+\frac{\tau\Delta}{2}}{\Delta}\right\rfloor & |x|\geq\tau\\
0 & \text{otherwise}
\end{cases}
\qquad
\hat{x}=
\begin{cases}
\operatorname{sign}(i)\,\Delta\left(|i|+\dfrac{1-\tau}{2}\right) & |x|\geq\tau\\
0 & \text{otherwise}
\end{cases}
\]

Parameter $\tau$ controls the dead-zone width: larger $\tau$ sends more samples to zero.

\begin{calloutbox}{red!70!black}{Why dead-zone helps lossy compression}
Natural signals are not sparse in their original domain, but become sparse after \textbf{prediction} or \textbf{linear transform}. A signal is sparse when most elements are zero or near zero. Setting near-zero values to zero adds little distortion, so the dead-zone produces many zero indexes. Sparse zero-heavy output is then cheap to entropy-code, which is where the real rate saving comes from.
\end{calloutbox}

This connects to the predictive and entropy-coding sections: a midrise quantizer maps even small fluctuations to a non-zero level, while a midtread dead-zone suppresses them, favoring sparsity.

\subsection{Granular Noise vs. Overload Noise}

\begin{itemize}
  \item \textbf{Granular noise}: input lies inside dynamic range, so $|e|\leq\Delta/2$.
  \item \textbf{Overload noise}: input exceeds dynamic range, so clipping creates potentially large error.
\end{itemize}

Design tension:

\begin{itemize}
  \item Smaller $\Delta$ gives lower granular noise but requires either more levels or narrower dynamic range.
  \item Smaller dynamic range increases clipping risk.
  \item Larger dynamic range avoids clipping but wastes levels and increases $\Delta$ for fixed $R$.
\end{itemize}

\subsection{Rate-Distortion Behavior}

For a uniform random variable and UQ:

\[
D=\sigma_X^2 2^{-2R}
\]

so every extra bit reduces distortion by about a factor 4 and increases SNR by about \textbf{6 dB}.

For generic sources at high resolution:

\[
D=K_X\sigma_X^2 2^{-2R}
\]

where $K_X$ depends on dynamic-range loading.

For high-resolution optimal scalar quantization:

\[
D=c_X\sigma_X^2 2^{-2R}
\]

where $c_X$ depends only on source PDF shape.

\subsection{Optimal Scalar Quantization}

Uniform quantizers are not always optimal. \textbf{Non-uniform scalar quantizers} allocate smaller cells where the input PDF is dense and larger cells where samples are rare.

The finite-rate optimum is found numerically with \textbf{Lloyd-Max}:

\begin{enumerate}
  \item Given reconstruction levels, choose thresholds by nearest neighbor.
  \item Given thresholds, choose reconstruction levels as centroids.
  \item Repeat until distortion stops decreasing.
\end{enumerate}

This is the scalar case of the same alternating idea behind \textit{k-means}.

\subsection{Predictive Scalar Quantization}

Scalar quantization ignores correlation. Predictive quantization uses past samples to predict the current one:

\[
y(n)=x(n)-v(n)
\]

where $v(n)$ is the prediction. If prediction is good, $\sigma_Y^2\ll\sigma_X^2$, so quantizing $y(n)$ is easier than quantizing $x(n)$ directly.

\begin{calloutbox}{red!70!black}{Sparse residual idea}
Natural images and audio are not sparse in their sample domain, but prediction can make residuals sparse. Sparse residuals produce many zero or near-zero quantizer indexes, which entropy coding can compress well.
\end{calloutbox}

\subsection{Drift Problem}

In predictive coding, encoder and decoder must predict from the \textbf{same reconstructed past samples}. If the encoder predicts from original samples but the decoder predicts from reconstructed samples, their predictors diverge and reconstruction drifts.

Correct predictive quantization is therefore \textbf{closed-loop}:

\begin{itemize}
  \item Encoder quantizes the residual.
  \item Encoder reconstructs locally.
  \item Encoder and decoder both update prediction buffers with reconstructed samples $\tilde{x}$.
\end{itemize}

\subsection{Entropy Coding Impact}

Using $R=\log_2 L$ assumes fixed-length indexes and hides the advantage of predictive quantization. Prediction errors are concentrated near zero, so variable-length or arithmetic coding assigns short codes to frequent zero indexes and longer codes to rare large errors.

Without entropy coding, predictive quantization may look disappointing. With entropy coding, it can provide large gains.

\section{Theory and Formulas}

\subsection{UQ of a Uniform Random Variable}

Let:

\[
X\sim\mathcal{U}\left(-\frac{A}{2},\frac{A}{2}\right)
\]

Then:

\[
f_X(x)=\frac{1}{A}, \qquad \sigma_X^2=\frac{A^2}{12}
\]

For UQ with $L$ cells:

\[
\sigma_Q^2
=\mathbb{E}[(X-\tilde{X})^2]
=\frac{1}{A}\sum_{i=1}^{L}\int_{\hat{x}^i-\Delta/2}^{\hat{x}^i+\Delta/2}(u-\hat{x}^i)^2\,du
\]

Each cell contributes $\Delta^3/12$, so:

\[
D=\sigma_Q^2=\frac{1}{A}L\frac{\Delta^3}{12}
=\frac{\Delta^2}{12}
\]

Since $\Delta=A/L=A2^{-R}$:

\[
D=\frac{A^2}{12L^2}
=\sigma_X^2 2^{-2R}
\]

\begin{calloutbox}{red!70!black}{6 dB per bit rule}
\[
\mathrm{SNR}
=10\log_{10}\frac{\sigma_X^2}{\sigma_X^2 2^{-2R}}
=10\log_{10}(2^{2R})
\approx 6R \quad [\mathrm{dB}]
\]

Each additional bit/sample improves SNR by about 6 dB under uniform high-resolution assumptions.
\end{calloutbox}

\subsection{High-Resolution UQ for a Generic Source}

High resolution means $L\to+\infty$ and each quantization cell is small. Locally, $p_X(x)$ is approximately constant, so the quantization error behaves like:

\[
e\sim\mathcal{U}\left(-\frac{\Delta}{2},\frac{\Delta}{2}\right)
\]

Thus:

\[
D=\frac{\Delta^2}{12}
=\frac{A^2}{12}2^{-2R}
=K_X\sigma_X^2 2^{-2R}
\]

with load factor:

\[
\gamma^2=\frac{A^2/4}{\sigma_X^2}, \qquad K_X=\frac{\gamma^2}{3}
\]

So:

\[
\mathrm{SNR}\approx 6R-10\log_{10}\frac{\gamma^2}{3}
\]

\textbf{Meaning:} poor dynamic-range loading costs SNR.

\subsection{High-Resolution Optimal Scalar Quantizer}

For a generic PDF and fixed rate:

\[
\sigma_Q^2=c_X\sigma_X^2 2^{-2R}
\]

where $U=X/\sigma_X$ and:

\[
c_X=\frac{1}{12}\left[\int_{\mathbb{R}}p_U^{1/3}(t)\,dt\right]^3
\]

\begin{center}\fitwidth{%
\begin{tabular}{|l|l|l|}
\hline
\textbf{Distribution} & \textbf{Shape factor $c_X$} & \textbf{Meaning} \\ \hline
Uniform & $1$ & best case for UQ \\ \hline
Gaussian & $\frac{\sqrt{3}}{2}\pi\approx2.72$ & heavier tails increase distortion \\ \hline
\end{tabular}}
\end{center}

\subsection{Lloyd-Max Conditions}

\begin{calloutbox}{red!70!black}{Lloyd-Max necessary conditions}
\textbf{Nearest-neighbor thresholds:}

\[
k=\arg\min_n |x-\hat{x}^n| \Rightarrow Q(x)=\hat{x}^k
\]

\[
t^i=\frac{\hat{x}^i+\hat{x}^{i+1}}{2}
\]

\textbf{Centroid reconstruction levels:}

\[
\hat{x}^i=
\frac{\int_{t^{i-1}}^{t^i}u\,p_X(u)\,du}
{\int_{t^{i-1}}^{t^i}p_X(u)\,du}
=\mathbb{E}[X\mid X\in\Theta^i]
\]

Distortion never increases at each iteration, but convergence to the global optimum is not guaranteed.
\end{calloutbox}

For a training set $\mathcal{X}=\{u_1,\ldots,u_M\}$, replace integrals with cluster assignments:

\[
W_k^i=\{u_m\in\mathcal{X}: \|u_m-\hat{x}_k^i\|\leq\|u_m-\hat{x}_k^j\|,\forall j\neq i\}
\]

\[
\hat{x}_{k+1}^i=\frac{1}{|W_k^i|}\sum_{u_m\in W_k^i}u_m
\]

\subsection{Prediction Gain}

Quantizing the prediction error:

\[
y(n)=x(n)-v(n)
\]

reconstructs:

\[
\hat{x}(n)=v(n)+\hat{y}(n)
\]

Prediction-error quantization error equals final reconstruction error:

\[
q(n)=y(n)-\hat{y}(n)
=[x(n)-v(n)]-[\hat{x}(n)-v(n)]
=x(n)-\hat{x}(n)
\]

Total SNR decomposes into prediction gain plus quantization gain:

\[
\mathrm{SNR}_P
=10\log_{10}\frac{\sigma_X^2}{D}
=
\underbrace{10\log_{10}\frac{\sigma_X^2}{\sigma_Y^2}}_{G_P}
+
\underbrace{10\log_{10}\frac{\sigma_Y^2}{D}}_{G_Q}
\]

\begin{calloutbox}{red!70!black}{Positive prediction gain}
Prediction helps if and only if:

\[
\sigma_Y^2<\sigma_X^2
\]

Equivalently, $G_P>0$.
\end{calloutbox}

\subsection{Linear Prediction and Wiener-Hopf Solution}

Order-$P$ linear predictor:

\[
v(n)=-\sum_{i=1}^{P}a_i x_{n-i}
\]

Prediction error:

\[
y(n)=x(n)-v(n)=\sum_{i=0}^{P}a_i x_{n-i}, \qquad a_0=1
\]

Prediction-error variance:

\[
\sigma_Y^2=\sigma_X^2+2\underline{r}^{T}\underline{a}
+\underline{a}^{T}R_X\underline{a}
\]

where:

\[
\underline{r}=[r_X(1),\ldots,r_X(P)]^T
\]

\[
(R_X)_{ij}=r_X(|i-j|), \qquad r_X(k)=\mathbb{E}[X(n)X(n-k)]
\]

\begin{calloutbox}{red!70!black}{Optimal predictor}
\[
\frac{\partial\sigma_Y^2}{\partial\underline{a}}
=2\underline{r}+2R_X\underline{a}=0
\]

\[
\underline{a}^{\mathrm{opt}}=-R_X^{-1}\underline{r}
\]

\[
\sigma_{Y,\mathrm{opt}}^2=\sigma_X^2+\underline{r}^{T}\underline{a}^{\mathrm{opt}}
\]
\end{calloutbox}

Autocorrelation can be estimated from data:

\[
\hat{r}_X(k)=\frac{1}{N}\sum_{n=0}^{N-1-k}X(n)X(n+k)
\]

\subsection{Local Adaptation and Side Information}

Natural images are non-stationary. Local prediction filters can track local textures better, but filter coefficients must be transmitted as side information.

For block size $M\times M$, filter order $N$, and $B$ bits/coefficient:

\[
R_{\text{overhead}}=\frac{NB}{M^2}
\]

Small blocks adapt well but have high overhead. Large blocks reduce overhead but adapt poorly.

\section{Visual Schemes}

\subsection{Scalar Quantizer Structure}

\begin{center}\includegraphics[width=\linewidth]{ch02/scalar-quantizer-model.png}\end{center}
Scalar quantization partitions the real line into regions and maps each region to one reconstruction level.

\begin{center}\includegraphics[width=0.633\linewidth]{ch02/midtread-quantizer.png}\end{center}
Midtread quantizer with odd $L$: zero is a reconstruction level.

\begin{center}\includegraphics[width=0.633\linewidth]{ch02/midrise-quantizer.png}\end{center}
Midrise quantizer with even $L$: zero is a threshold, not a reconstruction level.

\subsection{Rate-Distortion and Visual Artifacts}

\begin{center}\includegraphics[width=\linewidth]{ch02/grayscale-quantization-levels.png}\end{center}
Coarse scalar quantization creates visible banding and false contours, especially at low bit depth.

\begin{center}\includegraphics[width=\linewidth]{ch02/rate-distortion-curve.png}\end{center}
Rate-distortion curves show the expected reduction of distortion as bit rate increases.

\subsection{Predictive Coding Schemes}

\begin{center}\includegraphics[width=\linewidth]{ch02/open-loop-predictive-coding.png}\end{center}
Open-loop prediction quantizes the prediction error, but can drift if encoder and decoder predictors use different past samples.

\begin{center}\footnotesize\itshape Diagram (Mermaid source):\end{center}
\begin{verbatim}
flowchart LR
    X[Input x(n)] --> S[Subtract prediction]
    P[Predictor] --> S
    S --> Y[Residual y(n)]
    Y --> Q[Quantizer]
    Q --> I[Index stream]
    Q --> YH[Quantized residual]
    YH --> A[Add prediction]
    P --> A
    A --> XR[Local reconstruction]
    XR --> P
\end{verbatim}

\begin{center}\includegraphics[width=\linewidth]{ch02/prediction-residual-image.png}\end{center}
Prediction removes much of the local image structure, leaving a lower-variance residual.

\begin{center}\includegraphics[width=\linewidth]{ch02/predictor-order-effect.png}\end{center}
Higher predictor order improves residual whitening, but gains quickly saturate because near neighbors already explain most correlation.

\subsection{Drift and Correct Closed Loop}

\begin{center}\includegraphics[width=\linewidth]{ch02/predictive-coding-drift.png}\end{center}
Wrong scheme: encoder predicts from original samples, decoder predicts from reconstructed samples, so mismatch accumulates as drift.

\begin{center}\includegraphics[width=\linewidth]{ch02/predictive-coding-closed-loop.png}\end{center}
Correct scheme: encoder and decoder both predict from reconstructed samples, keeping state synchronized.

\begin{center}\footnotesize\itshape Diagram (Mermaid source):\end{center}
\begin{verbatim}
flowchart LR
    X[Input x(n)] --> M[Subtract]
    PR[Predictor from reconstructed past] --> M
    M --> Q[Quantize residual]
    Q --> EC[Entropy coding]
    EC --> BS[Bitstream]
    Q --> IQ[Decoded residual]
    IQ --> ADD[Add prediction]
    PR --> ADD
    ADD --> XH[Reconstructed sample]
    XH --> BUF[Past reconstructed buffer]
    BUF --> PR
\end{verbatim}

\subsection{Entropy Coding Gain}

\begin{center}\includegraphics[width=\linewidth]{ch02/prediction-error-index-distribution.png}\end{center}
Prediction errors produce many zero indexes, making fixed-length rate estimates too pessimistic.

\begin{center}\includegraphics[width=\linewidth]{ch02/predictive-quantization-entropy-gain.png}\end{center}
Predictive quantization becomes strongly effective when its sparse indexes are entropy-coded.

\section{Examples}

\begin{calloutbox}{green!55!black}{Image quantization at low rates}
\textbf{RGB peppers example:}

\begin{center}\fitwidth{%
\begin{tabular}{|l|l|l|}
\hline
\textbf{Rate} & \textbf{PSNR} & \textbf{Compression ratio} \\ \hline
8 bpp & 29.26 dB & 2.667 \\ \hline
6 bpp & 27.83 dB & 4.000 \\ \hline
3 bpp & 25.75 dB & 8.000 \\ \hline
\end{tabular}}
\end{center}

\textbf{Takeaway:} lower bit depth reduces rate, but creates banding and false contours.
\end{calloutbox}

\begin{calloutbox}{green!55!black}{Grayscale image quantization}
\begin{center}\fitwidth{%
\begin{tabular}{|l|l|l|l|}
\hline
\textbf{Quantization levels} & \textbf{Bits/pixel} & \textbf{Compression ratio} & \textbf{Y-PSNR} \\ \hline
256 & 8 & 1.000 & $\infty$ \\ \hline
64 & 6 & 1.333 & 46.51 dB \\ \hline
16 & 4 & 2.000 & 34.47 dB \\ \hline
8 & 3 & 2.667 & 29.06 dB \\ \hline
4 & 2 & 4.000 & 22.85 dB \\ \hline
2 & 1 & 8.000 & 17.42 dB \\ \hline
\end{tabular}}
\end{center}

\textbf{Takeaway:} rate reduction is smooth numerically, but visual quality drops sharply at low levels.
\end{calloutbox}

\begin{calloutbox}{green!55!black}{Audio quantization}
\begin{center}\fitwidth{%
\begin{tabular}{|l|l|}
\hline
\textbf{Bit depth} & \textbf{Perceptual quality} \\ \hline
8 bits & slight hissing or granular noise \\ \hline
4 bits & strong fidelity loss and prominent noise \\ \hline
2 bits & very strong noise; speech may remain intelligible, music is badly damaged \\ \hline
\end{tabular}}
\end{center}

\textbf{Takeaway:} scalar quantization noise is more tolerable for some signals than others, but very low bit depth destroys quality.
\end{calloutbox}

\begin{calloutbox}{green!55!black}{Prediction gain for AR(1)}
\textbf{Signal:}

\[
X(n)\sim\mathcal{N}(0,\sigma^2), \qquad
\mathbb{E}[X(n)X(m)]=\sigma^2\rho^{|n-m|}
\]

\textbf{Predictor:} $V(n)=X(n-1)$.

Prediction error variance:

\[
\sigma_Y^2
=\mathbb{E}[(X(n)-X(n-1))^2]
=2\sigma^2(1-\rho)
\]

Prediction gain:

\[
G_P=10\log_{10}\frac{\sigma^2}{2(1-\rho)\sigma^2}
=10\log_{10}\frac{1}{2(1-\rho)}
\]

Positive gain condition:

\[
G_P>0 \Leftrightarrow \rho>\frac{1}{2}
\]

For $\rho=0.9$, $G_P\approx7$ dB.
\end{calloutbox}

\begin{calloutbox}{green!55!black}{2D image predictor}
Predictor for grayscale image:

\[
\tilde{f}_{n,m}=a f_{n-1,m}+b f_{n,m-1}
\]

\begin{center}\fitwidth{%
\begin{tabular}{|l|l|l|l|}
\hline
\textbf{$a$} & \textbf{$b$} & \textbf{$\sigma_E^2$} & \textbf{Meaning} \\ \hline
0 & 0 & 2902.7 & no prediction \\ \hline
1/2 & 1/2 & 78.7 & simple neighbor average \\ \hline
0.449 & 0.546 & 78.4 & optimal predictor \\ \hline
\end{tabular}}
\end{center}

\textbf{Takeaway:} a simple local predictor reduces variance by more than 37 times; optimization adds only small improvement.
\end{calloutbox}

\begin{calloutbox}{green!55!black}{Drift with wrong prediction loop}
Encoder predicts from original past values, decoder predicts from reconstructed past values.

\begin{center}\fitwidth{%
\begin{tabular}{|l|l|l|l|l|l|l|}
\hline
\textbf{Step} & \textbf{Original} & \textbf{Encoder prediction} & \textbf{Encoder error} & \textbf{Quantized error} & \textbf{Decoder prediction} & \textbf{Reconstruction} \\ \hline
1 & 10 & - & - & 9 & - & 9 \\ \hline
2 & 11 & 10 & 1 & 0 & 9 & 9 \\ \hline
3 & 12 & 11 & 1 & 0 & 9 & 9 \\ \hline
4 & 13 & 12 & 1 & 0 & 9 & 9 \\ \hline
5 & 14 & 13 & 1 & 0 & 9 & 9 \\ \hline
6 & 18 & 14 & 4 & 3 & 9 & 12 \\ \hline
\end{tabular}}
\end{center}

\textbf{Takeaway:} decoder stays behind because it cannot reproduce encoder predictions.
\end{calloutbox}

\begin{calloutbox}{green!55!black}{No drift with closed-loop prediction}
Encoder and decoder both predict from reconstructed past values.

\begin{center}\fitwidth{%
\begin{tabular}{|l|l|l|l|l|l|l|}
\hline
\textbf{Step} & \textbf{Original} & \textbf{Encoder prediction} & \textbf{Encoder error} & \textbf{Quantized error} & \textbf{Decoder prediction} & \textbf{Reconstruction} \\ \hline
1 & 10 & - & - & 9 & - & 9 \\ \hline
2 & 11 & 9 & 2 & 3 & 9 & 12 \\ \hline
3 & 12 & 12 & 0 & 0 & 12 & 12 \\ \hline
4 & 13 & 12 & 1 & 0 & 12 & 12 \\ \hline
5 & 14 & 12 & 2 & 3 & 12 & 15 \\ \hline
6 & 18 & 15 & 3 & 3 & 15 & 18 \\ \hline
7 & 21 & 18 & 3 & 3 & 18 & 21 \\ \hline
8 & 18 & 21 & -3 & -3 & 21 & 18 \\ \hline
\end{tabular}}
\end{center}

\textbf{Takeaway:} shared reconstructed state keeps encoder and decoder synchronized.
\end{calloutbox}

\begin{calloutbox}{green!55!black}{Entropy coding gain}
With $L=19$ quantization levels, about \textbf{84\%} of prediction-error indexes are zero, and less than \textbf{1\%} fall outside $(-3,3)$.

\textbf{Result from source:} predictive quantization plus entropy coding gives about \textbf{+20 dB PSNR at 1 bpp} versus direct UQ, or about \textbf{86\% rate reduction} at 30 dB.

\textbf{Takeaway:} prediction creates sparsity; entropy coding turns sparsity into bitrate reduction.
\end{calloutbox}
